\section{Conclusion}
\label{sec:Con}

\par This study examined whether CVaR-minimisation provides stronger downside-risk control than Equal-Weight and Minimum-Variance portfolio construction in a Malta-accessible UCITS ETF basket. The research was structured around three questions covering the dataset, the portfolio-construction process, and the evaluation of out-of-sample performance.

\par RQ1 concerned the dataset and investable basket. The study showed that a fixed six-ETF UCITS basket could be used to create a consistent and transparent dataset for portfolio backtesting. The selected ETFs provided exposure to major equity markets and government-bond sleeves, while EUR alignment allowed the strategies to be compared under a common reporting currency. This provided a suitable controlled dataset for the scope of the research, although it remained smaller than the broader asset universes used in some portfolio-optimisation studies.

\par RQ2 concerned the portfolio-construction process. Under identical long-only and fully invested constraints, CVaR-minimisation behaved as a defensive allocation rule. It did not behave like Equal-Weight, which maintained constant exposure to all six ETFs, but was closer to Minimum-Variance in its preference for lower-risk and lower-correlation bond exposures. This confirmed that changing the portfolio-construction rule affected the resulting allocation, which is consistent with the reviewed literature on risk-metric choice and portfolio behaviour \cite{munoz2026riskmetrics,teplova2023blcopula}.

\par RQ3 concerned evaluation and testing. The out-of-sample results showed that Equal-Weight achieved the strongest return and risk-adjusted performance, with the highest CAGR and Sharpe ratio. CVaR-95 achieved the lowest VaR and CVaR values, which supports the part of the hypothesis related to daily tail-loss control. However, it did not achieve the lowest maximum drawdown, and its CAGR and Sharpe ratio were materially lower than Equal-Weight. Minimum-Variance achieved the best maximum drawdown. Therefore, the main hypothesis was only partially supported: CVaR-minimisation improved the specific tail-risk metrics it targeted, but it did not dominate the benchmark strategies across all downside-risk and return-adjusted measures.

\par Overall, the study achieved its main aim to a reasonable extent. It implemented the three strategies under the same dataset, constraints, rolling backtest design, and evaluation metrics, and produced a clear empirical comparison. However, the results also showed that a simple CVaR-minimisation framework is not automatically superior. Its defensive allocation reduced short-horizon tail-risk, but this came at the cost of weaker participation in equity-market growth. This explains why the CVaR strategy reduced VaR and CVaR, but did not preserve competitive CAGR or Sharpe performance.

\par The findings should be interpreted within the limits of the research design. The study used a small six-ETF UCITS basket, a specific 2017--2025 realised out-of-sample period, market-price data, and a simplified backtest without a full transaction-cost or tax model. These choices supported transparency and feasibility, but they also limit generalisability. The results should therefore be read as evidence from a controlled Malta-accessible ETF comparison, not as a universal conclusion about all CVaR-based portfolios.

\par Future research should extend this work in several directions. First, a larger UCITS ETF universe could be tested to see whether CVaR-minimisation performs differently when more asset classes and regions are available. Second, future work should include transaction costs, bid-ask spreads, taxes, and fund fees to make the backtest closer to real investor implementation. Third, researchers could compare exchange market-price data with official NAV-based fund data where available. Fourth, more advanced CVaR frameworks, such as mean-CVaR, target-return CVaR, Black-Litterman mean-CVaR, or regime-based optimisation, could be tested against the simple CVaR-minimisation approach used here. Finally, future studies could examine whether CVaR performs better during specific stress periods rather than over the whole out-of-sample period. These extensions would help determine whether the weakness found in this study was due to the CVaR objective itself, the small ETF universe, or the chosen market period.
